La radioactivitat és un mitjà fiable per a calcular l'edat de les roques i minerals que contenen isòtops radioactius concrets. Aquest sistema de datació radiomètrica ens permet mesurar el temps geològic.

Un d'aquests mètodes es basa en la desintegració de l'isòtop ${}^{40}_{19}\mathrm{K}$ (potassi) en ${}^{40}_{18}\mathrm{Ar}$ (argó). El rellotge potassi-argó comença a funcionar quan els minerals que contenen potassi cristal·litzen a partir d'un magma o dins una roca. En aquest moment, els nous minerals contenen ${}^{40}_{19}\mathrm{K}$ i no contenen ${}^{40}_{18}\mathrm{Ar}$. A mesura que passa el temps, el ${}^{40}_{19}\mathrm{K}$ es desintegra i tots els àtoms de ${}^{40}_{18}\mathrm{Ar}$ que trobem en el mineral en un temps posterior a la formació provenen de la descomposició del ${}^{40}_{19}\mathrm{K}$.

\begin{apartat}{1}
Escriviu la reacció nuclear de l'emissió de partícules $\beta$ de l'isòtop ${}^{40}_{19}\mathrm{K}$.
\end{apartat}

\begin{apartat}{1}
En una roca s'han trobat 10,0~g de ${}^{40}_{19}\mathrm{K}$ i 10,0~g de ${}^{40}_{18}\mathrm{Ar}$. Quina quantitat de ${}^{40}_{19}\mathrm{K}$ hi haurà quan hauran transcorregut $5,00\times 10^{9}$ anys? Fent servir la datació radiomètrica basada en el potassi-argó, digueu quina edat té la roca. Considereu que el ${}^{40}_{19}\mathrm{K}$ es desintegra només en ${}^{40}_{18}\mathrm{Ar}$.
\end{apartat}

\blocetiquetat{Dada:}{Període de semidesintegració del ${}^{40}_{19}\mathrm{K}$, $t_{1/2} = 1,25\times 10^{9}\ \text{anys}$}
